% =========================================================
% LAB 0
% =========================================================
\section{LAB0: Giới thiệu và làm quen}

\subsection{Trình bày, vẽ sơ đồ mạch theo cách hiểu của học viên}
Theo các tài liệu hướng dẫn và nội dung đã thiết kế, sơ đồ mạch của bài thực hành bao gồm các thành phần sau:
\begin{itemize}
    \item \textbf{Khối 1:} (Học viên điền tên và chức năng của module chính).
    \item \textbf{Khối 2:} (Học viên điền tên và chức năng của module hỗ trợ).
\end{itemize}

\begin{figure}[H]
    \centering
    % \includegraphics[width=0.8\textwidth]{figures/lab0_circuit.png}
    \vspace{4cm} % Khoảng trống chừa sẵn
    \caption{Sơ đồ mạch kiến trúc cho LAB0}
\end{figure}

\subsection{Sơ đồ dataflow với code software}
Sơ đồ mô phỏng luồng dữ liệu (Dataflow) liên kết với thiết kế và phần mềm:
\begin{figure}[H]
    \centering
    % \includegraphics[width=0.8\textwidth]{figures/lab0_dataflow.png}
    \vspace{4cm}
    \caption{Sơ đồ dataflow LAB0}
\end{figure}

Đoạn mã code software minh họa:
\begin{lstlisting}[language=C, caption={Code software minh họa LAB0}]
// Học viên chèn mã nguồn C/C++ hoặc Verilog vào đây
int main() {
    // Initialization
    
    // Main execution loop
    
    return 0;
}
\end{lstlisting}

\subsection{Báo cáo kết quả Quartus synthesis implementation, resource, $f_{max}$}
Sau khi thực hiện tổng hợp (synthesis) trên phần mềm Quartus Prime, kết quả về tài nguyên sử dụng và tần số hoạt động tối đa được thống kê chi tiết như sau:

\begin{table}[H]
    \centering
    \begin{tabular}{|l|l|}
        \hline
        \textbf{Thông số} & \textbf{Giá trị báo cáo} \\
        \hline
        Logic utilization (ALMs / LEs) & ... \\
        \hline
        Total registers & ... \\
        \hline
        Total pins & ... \\
        \hline
        Total memory bits & ... \\
        \hline
        $f_{max}$ (Maximum Frequency) & ... MHz \\
        \hline
    \end{tabular}
    \caption{Kết quả synthesis Resource và $f_{max}$ của LAB0}
\end{table}

\subsection{Hình ảnh chạy test code trên console}
\textit{(Chèn hình ảnh kết quả thực thi trên console tại đây)}
\begin{figure}[H]
    \centering
    % \includegraphics[width=0.8\textwidth]{figures/lab0_console.png}
    \vspace{5cm}
    \caption{Kết quả chạy test code trên console (LAB0)}
\end{figure}


% =========================================================
% LAB 1
% =========================================================
\newpage
\section{LAB1: Thiết kế và thực thi hệ thống 1}

\subsection{Trình bày, vẽ sơ đồ mạch theo cách hiểu của học viên}
Sơ đồ mạch của bài thực hành được trình bày như hình bên dưới:

\begin{figure}[H]
    \centering
    % \includegraphics[width=0.8\textwidth]{figures/lab1_circuit.png}
    \vspace{4cm}
    \caption{Sơ đồ mạch kiến trúc cho LAB1}
\end{figure}

\subsection{Sơ đồ dataflow với code software}
Dòng chảy dữ liệu của hệ thống:
\begin{figure}[H]
    \centering
    % \includegraphics[width=0.8\textwidth]{figures/lab1_dataflow.png}
    \vspace{4cm}
    \caption{Sơ đồ dataflow LAB1}
\end{figure}

Đoạn mã code software minh họa:
\begin{lstlisting}[language=C, caption={Code software minh họa LAB1}]
// Code phần mềm cho LAB1
\end{lstlisting}

\subsection{Báo cáo kết quả Quartus synthesis implementation, resource, $f_{max}$}
Kết quả synthesis Resource và $f_{max}$ của LAB1:

\begin{table}[H]
    \centering
    \begin{tabular}{|l|l|}
        \hline
        \textbf{Thông số} & \textbf{Giá trị báo cáo} \\
        \hline
        Logic utilization (ALMs / LEs) & ... \\
        \hline
        Total registers & ... \\
        \hline
        Total pins & ... \\
        \hline
        Total memory bits & ... \\
        \hline
        $f_{max}$ (Maximum Frequency) & ... MHz \\
        \hline
    \end{tabular}
    \caption{Kết quả synthesis Resource và $f_{max}$ của LAB1}
\end{table}

\subsection{Hình ảnh chạy test code trên console}
\textit{(Chèn hình ảnh kết quả thực thi trên console tại đây)}
\begin{figure}[H]
    \centering
    % \includegraphics[width=0.8\textwidth]{figures/lab1_console.png}
    \vspace{5cm}
    \caption{Kết quả chạy test code trên console (LAB1)}
\end{figure}


% =========================================================
% LAB 2
% =========================================================
\newpage
\section{LAB2: Thiết kế và thực thi hệ thống 2}

\subsection{Trình bày, vẽ sơ đồ mạch theo cách hiểu của học viên}
Sơ đồ mạch của bài thực hành được trình bày như hình bên dưới:

\begin{figure}[H]
    \centering
    % \includegraphics[width=0.8\textwidth]{figures/lab2_circuit.png}
    \vspace{4cm}
    \caption{Sơ đồ mạch kiến trúc cho LAB2}
\end{figure}

\subsection{Sơ đồ dataflow với code software}
Dòng chảy dữ liệu của hệ thống:
\begin{figure}[H]
    \centering
    % \includegraphics[width=0.8\textwidth]{figures/lab2_dataflow.png}
    \vspace{4cm}
    \caption{Sơ đồ dataflow LAB2}
\end{figure}

Đoạn mã code software minh họa:
\begin{lstlisting}[language=C, caption={Code software minh họa LAB2}]
// Code phần mềm cho LAB2
\end{lstlisting}

\subsection{Báo cáo kết quả Quartus synthesis implementation, resource, $f_{max}$}
Kết quả synthesis Resource và $f_{max}$ của LAB2:

\begin{table}[H]
    \centering
    \begin{tabular}{|l|l|}
        \hline
        \textbf{Thông số} & \textbf{Giá trị báo cáo} \\
        \hline
        Logic utilization (ALMs / LEs) & ... \\
        \hline
        Total registers & ... \\
        \hline
        Total pins & ... \\
        \hline
        Total memory bits & ... \\
        \hline
        $f_{max}$ (Maximum Frequency) & ... MHz \\
        \hline
    \end{tabular}
    \caption{Kết quả synthesis Resource và $f_{max}$ của LAB2}
\end{table}

\subsection{Hình ảnh chạy test code trên console}
\textit{(Chèn hình ảnh kết quả thực thi trên console tại đây)}
\begin{figure}[H]
    \centering
    % \includegraphics[width=0.8\textwidth]{figures/lab2_console.png}
    \vspace{5cm}
    \caption{Kết quả chạy test code trên console (LAB2)}
\end{figure}


% =========================================================
% LAB 3
% =========================================================
\newpage
\section{LAB3: Thiết kế và thực thi hệ thống 3}

\subsection{Trình bày, vẽ sơ đồ mạch theo cách hiểu của học viên}
Sơ đồ mạch của bài thực hành được trình bày như hình bên dưới:

\begin{figure}[H]
    \centering
    % \includegraphics[width=0.8\textwidth]{figures/lab3_circuit.png}
    \vspace{4cm}
    \caption{Sơ đồ mạch kiến trúc cho LAB3}
\end{figure}

\subsection{Sơ đồ dataflow với code software}
Dòng chảy dữ liệu của hệ thống:
\begin{figure}[H]
    \centering
    % \includegraphics[width=0.8\textwidth]{figures/lab3_dataflow.png}
    \vspace{4cm}
    \caption{Sơ đồ dataflow LAB3}
\end{figure}

Đoạn mã code software minh họa:
\begin{lstlisting}[language=C, caption={Code software minh họa LAB3}]
// Code phần mềm cho LAB3
\end{lstlisting}

\subsection{Báo cáo kết quả Quartus synthesis implementation, resource, $f_{max}$}
Kết quả synthesis Resource và $f_{max}$ của LAB3:

\begin{table}[H]
    \centering
    \begin{tabular}{|l|l|}
        \hline
        \textbf{Thông số} & \textbf{Giá trị báo cáo} \\
        \hline
        Logic utilization (ALMs / LEs) & ... \\
        \hline
        Total registers & ... \\
        \hline
        Total pins & ... \\
        \hline
        Total memory bits & ... \\
        \hline
        $f_{max}$ (Maximum Frequency) & ... MHz \\
        \hline
    \end{tabular}
    \caption{Kết quả synthesis Resource và $f_{max}$ của LAB3}
\end{table}

\subsection{Hình ảnh chạy test code trên console}
\textit{(Chèn hình ảnh kết quả thực thi trên console tại đây)}
\begin{figure}[H]
    \centering
    % \includegraphics[width=0.8\textwidth]{figures/lab3_console.png}
    \vspace{5cm}
    \caption{Kết quả chạy test code trên console (LAB3)}
\end{figure}
